\documentclass[aspectratio=169, xcolor=dvipsnames]{beamer}

% --- Theme Setup ---
\usetheme{Madrid}
\usecolortheme{default}

% --- Packages ---
\usepackage{graphicx}
\usepackage{xcolor}
\usepackage{booktabs}
\usepackage{hyperref}
\usepackage{adjustbox}
\usepackage{amssymb}
\usepackage{amsmath}
\usepackage{listings}
\usepackage{caption}
\usepackage{tikz}

\lstset{
  basicstyle=\ttfamily\small,
  keywordstyle=\color{blue}\bfseries,
  stringstyle=\color{red},
  showstringspaces=false,
  breaklines=true
}

% --- Title Information ---
\title{Module 50}
\subtitle{Concurrency Control/2}
\author{Partha Pratim Das}
\institute{Department of Computer Science and Engineering\\Indian Institute of Technology, Kharagpur\\ppd@cse.iitkgp.ac.in}
\date{\today}

\begin{document}

% Title Slide
\begin{frame}
    \titlepage
\end{frame}

\begin{frame}{Module Recap}
    \begin{itemize}[<+->]
        \item Understood the locking mechanism and protocols
        \item Realized that deadlock is a peril of locking and needs to be handled through rollback
    \end{itemize}
\end{frame}

\begin{frame}{Module Objectives}
    \begin{itemize}[<+->]
        \item Deadlocks are perils of locking. We need to understand how to detect, prevent and recover from deadlock
        \item Introduce a simple time-based protocol that avoids deadlocks
    \end{itemize}
\end{frame}

\begin{frame}{Module Outline}
    \begin{itemize}[<+->]
        \item Deadlock Handling
        \item Timestamp-Based Protocols
    \end{itemize}
\end{frame}

\section{Deadlock Handling}

\begin{frame}
  \vfill
  \centering
  \begin{beamercolorbox}[sep=8pt,center,shadow=true,rounded=true]{title}
    \usebeamerfont{title}Deadlock Handling\par%
  \end{beamercolorbox}
  \vfill
\end{frame}

\begin{frame}{Deadlock Handling}
    \begin{itemize}[<+->]
        \item System is \textbf{deadlocked} if there is a set of transactions such that every transaction in the set is waiting for another transaction in the set
        \item \textbf{Deadlock Prevention} protocols ensure that the system will never enter into a deadlock state. Some prevention strategies:
        \begin{itemize}
            \item[$\triangleright$] Require that each transaction locks all its data items before it begins execution (pre-declaration)
            \item[$\triangleright$] Impose partial ordering of all data items and require that a transaction can lock data items only in the order specified by the partial order
        \end{itemize}
    \end{itemize}
\end{frame}

\begin{frame}{Deadlock Prevention}
    \begin{block}{Transaction Timestamp}
        Timestamp is a unique identifier created by the DBMS to identify the relative starting time of a transaction. Timestamping is a method of concurrency control in which each transaction is assigned a transaction timestamp.
    \end{block}
    \begin{itemize}[<+->]
        \item Following schemes use transaction timestamps for the sake of deadlock prevention alone:
        \item \textbf{wait-die scheme}: non-preemptive
        \begin{itemize}
            \item[$\triangleright$] Older transaction may wait for younger one to release data item. (older means smaller timestamp)
            \item[$\triangleright$] Younger transactions never wait for older ones; they are rolled back instead
            \item[$\triangleright$] A transaction may die several times before acquiring needed data item
        \end{itemize}
        \item \textbf{wound-wait scheme}: preemptive
        \begin{itemize}
            \item[$\triangleright$] Older transaction wounds (forces rollback) of younger transaction instead of waiting for it
            \item[$\triangleright$] Younger transactions may wait for older ones
            \item[$\triangleright$] May be fewer rollbacks than wait-die scheme
        \end{itemize}
    \end{itemize}
\end{frame}

\begin{frame}{Deadlock Prevention (2): Wait-Die Scheme}
    \begin{itemize}[<+->]
        \item It is a \textbf{non-preemptive} technique for deadlock prevention
        \item When transaction $T_n$ requests a data item currently held by $T_k$, $T_n$ is allowed to wait only if it has a timestamp smaller than that of $T_k$ (That is, $T_n$ is older than $T_k$), otherwise $T_n$ is killed ('die')
        \item If a transaction requests to lock a resource (data item), which is already held with a conflicting lock by another transaction, then one of the two possibilities may occur:
        \begin{itemize}
            \item[$\triangleright$] \textbf{Timestamp($T_n$) $<$ Timestamp($T_k$)} : $T_n$, which is requesting a conflicting lock, is older than $T_k$, then $T_n$ is allowed to 'wait' until the data-item is available.
            \item[$\triangleright$] \textbf{Timestamp($T_n$) $>$ Timestamp($T_k$)} : $T_n$ is younger than $T_k$, then $T_n$ is killed ('dies'). $T_n$ is restarted later with a random delay but with the same timestamp(n)
        \end{itemize}
        \item This scheme allows the older transaction to 'wait' but kills the younger one ('die')
    \end{itemize}
\end{frame}

\begin{frame}{Deadlock Prevention (2): Wait-Die Scheme Example}
    \begin{exampleblock}{Example}
        Suppose that transactions $T_5$, $T_{10}$, $T_{15}$ have time-stamps 5, 10 and 15 respectively.
        \begin{itemize}
            \item If $T_5$ requests a data item held by $T_{10}$ then $T_5$ will 'wait'
            \item If $T_{15}$ requests a data item held by $T_{10}$, then $T_{15}$ will be killed ('die')
        \end{itemize}
    \end{exampleblock}
\end{frame}

\begin{frame}{Deadlock Prevention (3): Wound-Wait Scheme}
    \begin{itemize}[<+->]
        \item It is a \textbf{preemptive} technique for deadlock prevention
        \item When transaction $T_n$ requests a data item currently held by $T_k$, $T_n$ is allowed to wait only if it has a timestamp larger than that of $T_k$, otherwise $T_k$ is killed (wounded by $T_n$)
        \item If a transaction requests to lock a resource (data item), which is already held with a conflicting lock by another transaction, then one of the two possibilities may occur:
        \begin{itemize}
            \item[$\triangleright$] \textbf{Timestamp($T_n$) $<$ Timestamp($T_k$)} : $T_n$ forces $T_k$ to be killed ('wounds'). $T_k$ is restarted later with a random delay but with the same timestamp(k)
            \item[$\triangleright$] \textbf{Timestamp($T_n$) $>$ Timestamp($T_k$)} : $T_n$ 'wait's until the resource is free
        \end{itemize}
        \item This scheme allows the younger transaction requesting a lock to 'wait' if the older transaction already holds a lock, but forces the younger one to be suspended ('wound') if the older transaction requests a lock on an item already held by the younger one
    \end{itemize}
\end{frame}

\begin{frame}{Deadlock Prevention (3): Wound-Wait Scheme Example}
    \begin{exampleblock}{Example}
        Suppose that transactions $T_5$, $T_{10}$, $T_{15}$ have time-stamps 5, 10 and 15 respectively.
        \begin{itemize}
            \item If $T_5$ requests a data item held by $T_{10}$, then it will be preempted from $T_{10}$ and $T_{10}$ will be suspended ('wounded')
            \item If $T_{15}$ requests a data item held by $T_{10}$, then $T_{15}$ will 'wait'
        \end{itemize}
    \end{exampleblock}
\end{frame}

\begin{frame}{Deadlock Prevention}
    \begin{itemize}[<+->]
        \item Both in wait-die and in wound-wait schemes, a rolled back transaction is restarted with its original timestamp. Older transactions thus have precedence over newer ones, and starvation is hence avoided
    \end{itemize}
    \uncover<2->{
        \begin{block}{Timeout-Based Schemes}
            \begin{itemize}[<+->]
                \item A transaction waits for a lock only for a specified amount of time. If the lock has not been granted within that time, the transaction is rolled back and restarted
                \item Thus, deadlocks are not possible
                \item Simple to implement; but starvation is possible. Also difficult to determine good value of the timeout interval
            \end{itemize}
        \end{block}
    }
\end{frame}

\begin{frame}{Deadlock Detection}
    \begin{itemize}[<+->]
        \item Deadlocks can be described as a \textbf{wait-for graph}, which consists of a pair $G = (V, E)$
        \begin{itemize}
            \item[$\triangleright$] $V$ is a set of vertices (all the transactions in the system)
            \item[$\triangleright$] $E$ is a set of edges; each element is an ordered pair $T_i \rightarrow T_j$
        \end{itemize}
        \item If $T_i \rightarrow T_j$ is in $E$, then there is a directed edge from $T_i$ to $T_j$, implying that $T_i$ is waiting for $T_j$ to release a data item
        \item When $T_i$ requests a data item currently being held by $T_j$, then the edge $T_i \rightarrow T_j$ is inserted in the wait-for graph. This edge is removed only when $T_j$ is no longer holding a data item needed by $T_i$
        \item The system is in a deadlock state if and only if the wait-for graph has a \textbf{cycle}
        \item Must invoke a deadlock-detection algorithm periodically to look for cycles
    \end{itemize}
\end{frame}

\begin{frame}{Deadlock Recovery}
    \begin{itemize}[<+->]
        \item When deadlock is detected:
        \item \textbf{Some transaction will have to be rolled back (made a victim)} to break deadlock. Select that transaction as victim that will incur minimum cost
        \item \textbf{Rollback} - determine how far to roll back transaction
        \begin{itemize}
            \item[$\triangleright$] \textbf{Total rollback}: Abort the transaction and then restart it
            \item[$\triangleright$] \textbf{Partial rollback}: More effective to roll back transaction only as far as necessary to break deadlock
        \end{itemize}
        \item \textbf{Starvation} happens if same transaction is always chosen as victim. Include the number of rollbacks in the cost factor to avoid starvation
    \end{itemize}
\end{frame}

\section{Timestamp-Based Protocols}

\begin{frame}
  \vfill
  \centering
  \begin{beamercolorbox}[sep=8pt,center,shadow=true,rounded=true]{title}
    \usebeamerfont{title}Timestamp-Based Protocols\par%
  \end{beamercolorbox}
  \vfill
\end{frame}

\begin{frame}{Timestamp-Based Protocols}
    \begin{itemize}[<+->]
        \item Each transaction is issued a timestamp when it enters the system. If an old transaction $T_i$ has time-stamp $TS(T_i)$, a new transaction $T_j$ is assigned time-stamp $TS(T_j)$ such that $TS(T_i) < TS(T_j)$.
        \item The protocol manages concurrent execution such that the time-stamps determine the serializability order
        \item In order to assure such behavior, the protocol maintains for each data $Q$ two timestamp values:
        \begin{itemize}
            \item[$\triangleright$] \textbf{W-timestamp(Q)} is the largest time-stamp of any transaction that executed \texttt{write(Q)} successfully
            \item[$\triangleright$] \textbf{R-timestamp(Q)} is the largest time-stamp of any transaction that executed \texttt{read(Q)} successfully
        \end{itemize}
    \end{itemize}
\end{frame}

\begin{frame}{Timestamp-Based Protocols (2)}
    \begin{itemize}[<+->]
        \item The timestamp ordering protocol ensures that any conflicting read and write operations are executed in timestamp order
        \item Suppose a transaction $T_i$ issues a \texttt{read(Q)}:
        \begin{itemize}
            \item[$\triangleright$] a) If $TS(T_i) \le \text{W-timestamp}(Q)$, then $T_i$ needs to read a value of $Q$ that was already overwritten. Hence, the read operation is rejected, and $T_i$ is rolled back.
            \item[$\triangleright$] b) If $TS(T_i) \ge \text{W-timestamp}(Q)$, then the read operation is executed, and \text{R-timestamp}(Q) is set to $\max(\text{R-timestamp}(Q), TS(T_i))$.
        \end{itemize}
    \end{itemize}
\end{frame}

\begin{frame}{Timestamp-Based Protocols (3)}
    \begin{itemize}[<+->]
        \item Suppose that transaction $T_i$ issues \texttt{write(Q)}:
        \begin{itemize}
            \item[$\triangleright$] If $TS(T_i) < \text{R-timestamp}(Q)$, then the value of $Q$ that $T_i$ is producing was needed previously, and the system assumed that that value would never be produced. Hence, the write operation is rejected, and $T_i$ is rolled back
            \item[$\triangleright$] If $TS(T_i) < \text{W-timestamp}(Q)$, then $T_i$ is attempting to write an obsolete value of $Q$. Hence, this write operation is rejected, and $T_i$ is rolled back
            \item[$\triangleright$] Otherwise, the write operation is executed, and \text{W-timestamp}(Q) is set to $TS(T_i)$
        \end{itemize}
    \end{itemize}
\end{frame}

\begin{frame}{Correctness of Timestamp-Ordering Protocol}
    \begin{itemize}[<+->]
        \item The timestamp-ordering protocol guarantees serializability since all the arcs in the precedence graph are of the form:
        \begin{itemize}
            \item[$\triangleright$] $T_i \rightarrow T_j$ implies $TS(T_i) < TS(T_j)$
        \end{itemize}
        \item Thus, there will be no cycles in the precedence graph
        \item Timestamp protocol ensures freedom from deadlock as no transaction ever waits
        \item But the schedule may not be cascade-free, and may not even be recoverable
    \end{itemize}
\end{frame}

\begin{frame}{Module Summary}
    \begin{itemize}[<+->]
        \item Explained how to detect, prevent and recover from deadlock
        \item Introduced a time-based protocol that avoids deadlocks
    \end{itemize}
    \vspace{2em}
    \uncover<3->{
        \begin{center}
        \small \textit{Slides used in this presentation are borrowed from \url{http://db-book.com/} with kind permission of the authors.} \\
        \small \textit{Edited and new slides are marked with 'PPD'.}
        \end{center}
    }
\end{frame}

\end{document}
